\documentclass[dvisvgm,tikz,border=3pt]{standalone}
\usepackage{amsmath,amssymb}
\usepackage{pgfplots}
\usepackage{CJKutf8}
\pgfplotsset{compat=1.18}
\usetikzlibrary{arrows.meta,calc,positioning}

\definecolor{themebg}{HTML}{2e362c}
\definecolor{themefg}{HTML}{edf4e8}
\definecolor{themegray}{HTML}{9ea897}
\definecolor{themecurve}{HTML}{7eb6ff}
\definecolor{themealert}{HTML}{ff7b7b}
\definecolor{themeamber}{HTML}{f5b942}

\pgfdeclarelayer{background}
\pgfdeclarelayer{foreground}
\pgfsetlayers{background,main,foreground}

\begin{document}
\begin{CJK*}{UTF8}{gbsn}
\pagecolor{themebg}
\begin{tikzpicture}[color=themefg, text=themefg]

% ── 上图: 自然数下标轴抽样与覆盖失败 ──
\begin{axis}[
    name=ax1,
    width=13.8cm,
    height=4.8cm,
    axis lines=middle,
    axis line style={color=themefg, -{Stealth}},
    tick style={color=themefg},
    tick label style={color=themefg, font=\scriptsize},
    every axis label/.style={color=themefg},
    xmin=-0.5, xmax=11.0,
    ymin=-0.6, ymax=1.8,
    xtick={0,1,2,3,4,5,6,7,8,9},
    ytick=\empty,
    clip=false,
    xlabel={$n$},
    xlabel style={at={(1.0,0.18)}, anchor=south},
    title style={font=\footnotesize\bfseries, at={(0.5,1.08)}, anchor=south},
    title={子列抽样 vs 完整覆盖：$3n$ 与 $3n+1$ 遗漏全部 $3n+2$ 模类},
]

% 坐标轴上标出不同模类的点
% 3n (0, 3, 6, 9)
\addplot[only marks, mark=*, mark size=3pt, fill=themecurve, draw=themefg] coordinates {(0,0.9) (3,0.9) (6,0.9) (9,0.9)};
\node[text=themecurve, font=\scriptsize, fill=themebg, inner sep=0.8pt, anchor=south] at (axis cs:0,1.05) {$a_0$};
\node[text=themecurve, font=\scriptsize, fill=themebg, inner sep=0.8pt, anchor=south] at (axis cs:3,1.05) {$a_3$};
\node[text=themecurve, font=\scriptsize, fill=themebg, inner sep=0.8pt, anchor=south] at (axis cs:6,1.05) {$a_6$};
\node[text=themecurve, font=\scriptsize, fill=themebg, inner sep=0.8pt, anchor=south] at (axis cs:9,1.05) {$a_9$};
\node[text=themecurve, font=\scriptsize, fill=themebg, inner sep=0.8pt, anchor=west] at (axis cs:9.6,1.35) {$\{3n\}\to 0$};

% 3n+1 (1, 4, 7)
\addplot[only marks, mark=square*, mark size=2.8pt, fill=themeamber, draw=themefg] coordinates {(1,0.45) (4,0.45) (7,0.45)};
\node[text=themeamber, font=\scriptsize, fill=themebg, inner sep=0.8pt, anchor=south] at (axis cs:1,0.6) {$a_1$};
\node[text=themeamber, font=\scriptsize, fill=themebg, inner sep=0.8pt, anchor=south] at (axis cs:4,0.6) {$a_4$};
\node[text=themeamber, font=\scriptsize, fill=themebg, inner sep=0.8pt, anchor=south] at (axis cs:7,0.6) {$a_7$};
\node[text=themeamber, font=\scriptsize, fill=themebg, inner sep=0.8pt, anchor=west] at (axis cs:9.6,0.9) {$\{3n+1\}\to 0$};

% 3n+2 遗漏 (2, 5, 8)
\addplot[only marks, mark=x, mark size=4pt, line width=1.5pt, themealert] coordinates {(2,0.0) (5,0.0) (8,0.0)};
\node[text=themealert, font=\scriptsize, fill=themebg, inner sep=0.8pt, anchor=north] at (axis cs:2,-0.12) {$a_2=1$};
\node[text=themealert, font=\scriptsize, fill=themebg, inner sep=0.8pt, anchor=north] at (axis cs:5,-0.12) {$a_5=1$};
\node[text=themealert, font=\scriptsize, fill=themebg, inner sep=0.8pt, anchor=north] at (axis cs:8,-0.12) {$a_8=1$};
\node[text=themealert, font=\scriptsize, fill=themebg, inner sep=0.8pt, anchor=west] at (axis cs:9.6,0.45) {遗漏 $\{3n+2\}\to 1$};

\end{axis}

% ── 下图: 奇偶两子列趋向不同终点 A != B 否定收敛 ──
\begin{axis}[
    name=ax2,
    at={($(ax1.south)-(0,1.4cm)$)},
    anchor=north,
    width=13.8cm,
    height=5.4cm,
    axis lines=middle,
    axis line style={color=themefg, -{Stealth}},
    tick style={color=themefg},
    tick label style={color=themefg, font=\scriptsize},
    every axis label/.style={color=themefg},
    xmin=-0.5, xmax=11.0,
    ymin=-1.8, ymax=1.8,
    xtick={1,2,3,4,5,6,7,8,9,10},
    ytick={-1,0,1},
    clip=false,
    xlabel={$n$},
    xlabel style={at={(1.0,0.45)}, anchor=south},
    ylabel={$a_n$},
    title style={font=\footnotesize\bfseries, at={(0.5,1.08)}, anchor=south},
    title={两子列终点冲突：$a_{2n}\to A$ 与 $a_{2n-1}\to B$ ($A\ne B$) 直接否定收敛},
]
\begin{pgfonlayer}{background}
    % 上极限线 y=1 (A)
    \draw[themecurve, dashed, line width=1pt] (axis cs:0,1) -- (axis cs:10.5,1);
    \node[text=themecurve, font=\scriptsize, fill=themebg, inner sep=0.8pt, anchor=south west] at (axis cs:0.1,1.05) {$A=1$};

    % 下极限线 y=-1 (B)
    \draw[themealert, dashed, line width=1pt] (axis cs:0,-1) -- (axis cs:10.5,-1);
    \node[text=themealert, font=\scriptsize, fill=themebg, inner sep=0.8pt, anchor=north west] at (axis cs:0.1,-1.05) {$B=-1$};
\end{pgfonlayer}

% 偶数项靠近 1: 1 + 1/n
\addplot[only marks, mark=*, mark size=2.5pt, fill=themecurve, draw=themefg] coordinates {
    (2, 1.5) (4, 1.25) (6, 1.166) (8, 1.125) (10, 1.1)
};
\node[text=themecurve, font=\scriptsize, fill=themebg, inner sep=0.8pt, anchor=south west] at (axis cs:2,1.55) {$a_2$};
\node[text=themecurve, font=\scriptsize, fill=themebg, inner sep=0.8pt, anchor=south west] at (axis cs:4,1.3) {$a_4$};

% 奇数项靠近 -1: -1 - 1/n
\addplot[only marks, mark=*, mark size=2.5pt, fill=themealert, draw=themefg] coordinates {
    (1, -1.5) (3, -1.333) (5, -1.2) (7, -1.142) (9, -1.111)
};
\node[text=themealert, font=\scriptsize, fill=themebg, inner sep=0.8pt, anchor=north west] at (axis cs:1,-1.55) {$a_1$};
\node[text=themealert, font=\scriptsize, fill=themebg, inner sep=0.8pt, anchor=north west] at (axis cs:3,-1.35) {$a_3$};

\end{axis}

% 底部总结文字 (无框、高对比度)
\node[text=themefg, font=\scriptsize, anchor=north] at ($(ax2.south)-(0,0.25cm)$) {\textbf{判别准则}：证明存在必须控制全部尾部下标；否定收敛只需一对聚点冲突 ($A\ne B$)。};

\end{tikzpicture}
\end{CJK*}
\end{document}
